\chapter{Telepítés, binárisok és futtatási topológiák}\label{ch:telepites}

\noindent\textbf{Tanulási út: Gyors út: alap.} Juss el ismételhető helyi check/build/run workflow-ig; első olvasáskor hagyd későbbre a backend belső részleteit.\par\medskip

\section{Build a forrásból}

A kiinduló release forráscsomag Rust workspace. A két telepíthető program előállításához a szokásos release build:

\begin{lstlisting}[language=bash]
cargo build --locked --release
\end{lstlisting}

A repository saját teljes önellenőrzése a \code{./verify.sh}; ezt akkor futtasd, amikor magát a Velran forrásfát validálod vagy release-eled. Egy alkalmazás használatához nem ez az első lépés.

A \code{--locked} biztosítja, hogy a build a release-hez tartozó \code{Cargo.lock} függőséggráfot használja. Ha csak a két futtatható program kell:

\begin{lstlisting}[language=bash]
cargo build --locked --release -p velran-server -p velran-cli
\end{lstlisting}

A build után a két fontos program:

\begin{lstlisting}
target/release/velran-server
target/release/velran-cli
\end{lstlisting}

\begin{itemize}
  \item \code{velran-server}: a Velran alkalmazás fordítását és HTTP(S) kiszolgálását végző hosszú életű service process;
  \item \code{velran-cli}: operátori/adminisztratív program, például auth- és migrációs műveletekhez.
\end{itemize}

A két bináris szerepe szándékosan különbözik: a reverse proxy kizárólag a \code{velran-server} processhez kapcsolódik. A \code{velran-cli} nem webserver és nem kell publikus vagy proxy mögötti listener.

\subsection*{Futtatási követelmény: a rustc a Velran runtime toolchain része}

Minden olyan gépen, amely Velran alkalmazást futtat, rendelkezésre kell állnia egy használható \code{rustc} fordítónak. A Velrannak nincs interpreter- vagy VM-fallbackje: az alkalmazás backendje az ellenőrzött Velran kódot safe Rust kóddá alakítja, majd a \code{rustc} állítja elő a szerver által aktivált immutable natív \code{cdylib}-et. Ezért önmagában a \code{velran-server} telepítése \code{rustc} nélkül nem elegendő sem alkalmazásaktiváláshoz, sem source reloadhoz.

Alapértelmezésben a szerver az aktív toolchainből ismeri fel a \code{rustc}-t. Production telepítésnél a \code{server.toml} \code{[native]} szekciójában explicit, abszolút compiler-útvonal is megadható. A Cargo magának a Velrannak a forrásból történő buildjéhez kell; az alkalmazás-végrehajtási backend közvetlenül a \code{rustc}-t hívja.

\subsection*{Amit az alkalmazás aktiválásáról kezdetben tudnod kell}

Normál webfejlesztéshez először csak ezt a szabályt jegyezd meg: a Velran egy teljes candidate alkalmazást ellenőriz és fordít le, mielőtt az leválthatná az aktuálisan aktív verziót. Ha a candidate hibás, az utolsó ismert jó verzió szolgál tovább. Nincs interpreter- vagy VM-fallback, amely a nem támogatott alkalmazáskódot csendben futtatná.

Az első alkalmazások elkészítéséhez nem kell ismerned a natív compiler pipeline részleteit. Ezek szándékosan később, az ``A motorháztető alatt'' részben és a production fejezetben jelennek meg.

Fejlesztés közben elegendő lehet:

\begin{lstlisting}[language=bash]
cargo check --workspace
cargo test --workspace
\end{lstlisting}

\section{A binárisok telepítési helye}

Production gépen ne a Cargo \code{target/} könyvtárából futtasd tartósan a szolgáltatást. Kézi/forrásos telepítésnél a könyv a \code{/usr/local/bin} alá teszi a binárisokat, a \code{/usr/local/etc/velran} alá a géphez tartozó konfigurációt; így ezek elkülönülnek a disztribúció által kezelt fájloktól.

Egy egyszerű, jól auditálható elrendezés:

\begin{lstlisting}
/usr/local/bin/velran-server
/usr/local/bin/velran-cli
/usr/local/etc/velran/server.toml
/srv/velran/current/
/srv/velran/data/
/run/secrets/velran/
/var/log/velran/
\end{lstlisting}

Telepítés például:

\begin{lstlisting}[language=bash]
sudo install -d -m 0755 /usr/local/bin
sudo install -d -m 0755 /usr/local/etc/velran
sudo install -m 0755 target/release/velran-server /usr/local/bin/velran-server
sudo install -m 0755 target/release/velran-cli    /usr/local/bin/velran-cli
sudo install -m 0644 config/server.toml.sample \\
  /usr/local/etc/velran/server.toml
\end{lstlisting}

A \code{/usr/local/bin} a szokásos rendszer-PATH része, ezért külön symlinkre nincs szükség. A könyv példáiban a valódi binárisneveket, \code{velran-server} és \code{velran-cli}, használjuk.


\subsection*{Debian csomag esetén szándékosan más az elrendezés}

A fenti elrendezés a kézi/forrásból telepítésre vonatkozik. Ha a Velran Debian csomagból kerül a gépre, a fájlokat a \code{dpkg} kezeli, ezért a Debian filesystem-konvenció érvényes, nem a lokális adminisztrátori \code{/usr/local} prefix:

\begin{lstlisting}
/usr/bin/velran-server
/usr/bin/velran-cli
/etc/velran/server.toml
/usr/lib/systemd/system/velran.service
\end{lstlisting}

A csomag Debian/Ubuntu build gépen így készíthető és telepíthető:

\begin{lstlisting}[language=bash]
make deb
sudo dpkg -i dist/velran_1.0.0-1_amd64.deb
\end{lstlisting}

A csomag szándékosan nem indítja el automatikusan a szolgáltatást: telepítsd az alkalmazást és a secreteket, állítsd be és validáld az \code{/etc/velran/server.toml} fájlt, majd engedélyezd az \code{velran.service} unitot. Csomagolt telepítésnél \code{/usr} és \code{/etc}, kézi telepítésnél \code{/usr/local} a helyes. A teljes folyamat: \code{docs/hu/36-debian-csomag.md}.

\section{Első alkalmazás}

Ajánlott minimális projekt:

\begin{lstlisting}
myapp/
|-- main.vrn
|-- pages.vrn
`-- public/
\end{lstlisting}

\code{main.vrn}:

\begin{lstlisting}[language=Velran]
mod pages;
\end{lstlisting}

\code{pages.vrn}:

\begin{lstlisting}[language=Velran]
#[page]
fn home(ctx: PageContext) -> Result<Html, PageError> {
    return Ok(html {
        <main><h1>Hello Velran</h1></main>
    });
}

route home GET "/" public => home;
\end{lstlisting}

Fejlesztői indítás közvetlenül a workspace-ből:

\begin{lstlisting}[language=bash]
cargo run -p velran-server -- \\
  --app /abs/path/myapp/main.vrn \\
  --listen 127.0.0.1:8080 \\
  --insecure-dev-cookies
\end{lstlisting}

Ugyanez a lefordított release binárissal:

\begin{lstlisting}[language=bash]
/usr/local/bin/velran-server \\
  --app /abs/path/myapp/main.vrn \\
  --listen 127.0.0.1:8080 \\
  --insecure-dev-cookies
\end{lstlisting}

Az \code{--insecure-dev-cookies} csak lokális developmentre való.

\section{A konfigurációból elsőre csak a minimum kell}

A könyv ezen pontján még nem cél a teljes \code{server.toml} megtanulása. Kezdőként elég három dolgot szétválasztani:

\begin{itemize}
  \item az \emph{alkalmazáskód} mondja meg, milyen modellek, queryk, oldalak és route-ok léteznek;
  \item a \emph{szerverkonfiguráció} mondja meg, hol hallgat a process, milyen külső erőforrásokat érhet el, és milyen operátori limitek érvényesek;
  \item a \emph{secret fájlok} tartalmazzák azokat az érzékeny értékeket, amelyek nem kerülhetnek sem az alkalmazáskódba, sem plaintext TOML-ba.
\end{itemize}

Az A--Z fejezet egy minimális development konfigurációt ad. A production-specifikus kulcsokat akkor vezetjük be, amikor már van mit üzemeltetni; a teljes kulcs/default/sample referencia a könyv végén található. Első olvasáskor tehát ne próbáld a teljes konfigurációs felületet memorizálni.

\subsection*{A motorháztető alatt: natív alkalmazásaktiválás}

Első olvasáskor ezt a részt nyugodtan átugorhatod. Belsőleg az aktiválási út tudatosan lépcsőzetes, hogy az ellenőrzés azelőtt történjen meg, hogy az új kód éles generációvá válhatna:

\begin{lstlisting}
.vrn forras
  -> parser + security/type verification
  -> verified executable intermediate representation (IR)
  -> generalt safe Rust
  -> rustc
  -> immutable cdylib generation
  -> atomic activation
\end{lstlisting}

A generált alkalmazás-Rust tiltja az unsafe kódot. A nem támogatott konstrukció fail-closed hibát ad interpreter/VM kerülőút helyett. Source reload közben a hibás candidate ezért nem váltja le az utolsó ismert jó generációt. A natív compiler cache a változatlan, ellenőrzött source/toolchain/policy identity újrafordítását kerüli el.

\section{Production process és reverse proxy}

Apache vagy Nginx mögött a tipikus adatút:

\begin{lstlisting}
Internet
   |
   | HTTPS :443
   v
Apache vagy Nginx
   |
   | HTTP 127.0.0.1:8080
   v
/usr/local/bin/velran-server
   |
   +-- /usr/local/etc/velran/server.toml
   +-- /srv/velran/current/main.vrn
   +-- DB / Redis / AppFs
\end{lstlisting}

A proxy nem indítja és nem helyettesíti a Velran processzt. A \code{velran-server} külön service-ként fusson, például systemd alatt, a proxy pedig a loopback listenerre továbbítsa a kéréseket. A \code{velran-cli} csak operátori parancsok idejére indul el.

Egyszerű systemd \code{ExecStart} részlet:

\begin{lstlisting}
ExecStart=/usr/local/bin/velran-server \\
  --config /usr/local/etc/velran/server.toml
\end{lstlisting}

A process felhasználója ne legyen \code{root}; csak a szükséges alkalmazás-, secret-, data- és log-fájlokhoz kapjon jogosultságot. Ha a TLS-t Apache vagy Nginx terminálja, a Velran listener maradjon loopbacken, például \code{127.0.0.1:8080}.

\section{Futtatási topológiák részletesen}

A binárisok helye és a processzek szerepe ugyanaz marad, akár a Velran terminálja a TLS-t, akár Apache vagy Nginx áll előtte. A különbség a listener, a TLS-termináció és a trusted proxy határ konfigurációja.

\section{Közvetlen, egydomaines HTTPS}

Ha a Velran szerver maga terminálja a TLS-t, a publikus listener lehet például 443:

\begin{lstlisting}
[server]
app = "/srv/velran/current/main.vrn"
listen = "0.0.0.0:443"
insecure_dev_cookies = false

[tls]
cert_file = "/run/secrets/velran/tls-fullchain.pem"
key_file = "/run/secrets/velran/tls-key.pem"
handshake_timeout_ms = 5000
http_redirect_listen = "0.0.0.0:80"
public_host = "www.example.com"
\end{lstlisting}

A \code{public_host} trusted konfigurációból jön. HTTP-ről HTTPS-re csak GET/HEAD kérés redirectelhető.

Mivel a közvetlen minta 80/443 portot használ, unprivileged systemd service esetén a processnek célzott \code{CAP_NET_BIND_SERVICE} capability kell. Ne futtasd emiatt rootként. Reverse proxy mögött, például \code{127.0.0.1:8080} listenerrel ez a capability nem szükséges.

\section{Apache reverse proxy mögött}

Ebben a felállásban az Apache kezeli a publikus TLS-t, a Velran pedig csak loopback interfészen hallgat. Példa Velran config:

\begin{lstlisting}
[server]
app = "/srv/velran/current/main.vrn"
listen = "127.0.0.1:8080"
insecure_dev_cookies = false

[tls]
public_host = "www.example.com"

[web]
trusted_proxy_cidrs = ["127.0.0.1/32", "::1/128"]
allow_missing_origin = false
cors_origins = []
cors_allow_credentials = false
\end{lstlisting}

Reverse proxy módban nincs Velran oldali certificate/key, de a \code{tls.public_host} továbbra is szükséges: ez a böngésző által látott canonical host trusted forrása, amelyhez a Host és az Origin/Referer ellenőrzés kötődik. A production plain HTTP upstream csak loopback listeneren, explicit trusted proxy listával és ilyen public hosttal engedett. Minden normál proxied kérésnek trusted proxytól származó \code{X-Forwarded-Proto: https} (vagy ekvivalens szabványos \code{Forwarded: proto=https}) metadata-t kell hordoznia.

Az Apache VirtualHost minimális mintája:

\begin{lstlisting}
<VirtualHost *:443>
    ServerName www.example.com

    SSLEngine on
    SSLCertificateFile /etc/letsencrypt/live/www.example.com/fullchain.pem
    SSLCertificateKeyFile /etc/letsencrypt/live/www.example.com/privkey.pem

    ProxyPreserveHost On
    ProxyPass        / http://127.0.0.1:8080/
    ProxyPassReverse / http://127.0.0.1:8080/

    RequestHeader set X-Forwarded-Proto "https"
</VirtualHost>
\end{lstlisting}

Ehhez tipikusan a \code{proxy}, \code{proxy_http}, \code{headers}, \code{ssl} modulok szükségesek. A forwarding headereket a Velran csak a \code{trusted_proxy_cidrs} tartományból fogadja el; ezért ne tegyél tág, például teljes internetet lefedő trusted proxy CIDR-t a configba.

A port 80-as VirtualHost feladata lehet kizárólag az Apache-s HTTPS redirect:

\begin{lstlisting}
<VirtualHost *:80>
    ServerName www.example.com
    Redirect permanent / https://www.example.com/
</VirtualHost>
\end{lstlisting}

\section{Nginx reverse proxy mögött}

Az Nginx topológia ugyanazt a trust modellt használja: a publikus TLS-t az Nginx kezeli, a Velran pedig loopbacken hallgat. A Velran oldali \code{server.toml} megegyezhet az Apache példában látottal.

Minimális Nginx server block:

\begin{lstlisting}
server {
    listen 443 ssl;
    server_name www.example.com;

    ssl_certificate     /etc/letsencrypt/live/www.example.com/fullchain.pem;
    ssl_certificate_key /etc/letsencrypt/live/www.example.com/privkey.pem;

    location / {
        proxy_pass http://127.0.0.1:8080;
        proxy_http_version 1.1;
        proxy_set_header Host $host;
        proxy_set_header X-Forwarded-Proto https;
        proxy_set_header X-Forwarded-For $remote_addr;
    }
}
\end{lstlisting}

A példában egyetlen közvetlen reverse proxy van, ezért az \code{X-Forwarded-For} értéke a proxy által látott klienscím, az \code{X-Forwarded-Proto} pedig \code{https}. Ne engedd, hogy az internet felől érkező tetszőleges forwarding header authority legyen. A Velran csak a pontos \code{trusted_proxy_cidrs} tartományból érkező forwarding headert fogadja el.

A HTTP listener feladata lehet csak a HTTPS redirect:

\begin{lstlisting}
server {
    listen 80;
    server_name www.example.com;
    return 301 https://$host$request_uri;
}
\end{lstlisting}

\section{Ajánlott production könyvtárstruktúra}

\begin{lstlisting}
/usr/local/bin/velran-server
/usr/local/bin/velran-cli
/srv/velran/releases/2026-09-04.1/
/srv/velran/current -> /srv/velran/releases/2026-09-04.1
/srv/velran/data/
/usr/local/etc/velran/server.toml
/run/secrets/velran/
/var/log/velran/
\end{lstlisting}

A bináris, a release és a konfiguráció legyen runtime read-only. Az alkalmazás csak a deklarált data/log helyekre írjon.

\section{Első ellenőrzés}

Mielőtt forgalmat adnál a szerverre:

\begin{lstlisting}[language=bash]
/usr/local/bin/velran-server \\
  --config /usr/local/etc/velran/server.toml \\
  --check-config
\end{lstlisting}

Ez listener nélkül ellenőrzi többek között a konfigurációt, a secret-file referenciákat, az alkalmazás fordíthatóságát és a TLS fájlokat.

\section{Gyakorlat: ellenőrizd a telepítési határokat}
\paragraph{Gyakorlási mód: vezetett.} Használd az előszóban meghatározott vezetett módszert.

Egy tesztgépen azonosítsd a binárist, az immutable release-fát, a konfigurációt, a secret fájlokat, az alkalmazásadatot és a logokat. Ellenőrizd, hogy a service account csak a szándékosan írható adat- és loghelyekre tud írni. Listener nyitása előtt futtasd a \code{--check-config} ellenőrzést.

\section{Tipikus hiba: a konfigurációt alkalmazásadatként kezeljük}
A konfiguráció és a secret ne sodródjon együtt a módosuló alkalmazásállapottal. A read-only release-fa és konfiguráció jóval egyszerűbbé teszi a rollbacket, a felülvizsgálatot és az incidenskezelést.

\section{Folyamatos mintaprojekt: Secure Notes}
A Secure Notes esetén döntsd el, hol él az alkalmazás release artefaktuma, hol van a tartós adatbázis vagy adatkönyvtár, és mely secretek érkeznek fájlból ahelyett, hogy forráskódba kerülnének. Ezeket az útvonalakat még a feature-ök megvalósítása előtt rögzítsd.
